La mémoire de l'ordinateur ne contient que des zéros et des uns. Il faut donc pouvoir 
représenter toutes les informations que l'on veut mémoriser par des suites de zéros et
de uns. Cette traduction de l'information en binaire s'appelle le {\em codage binaire}. Les chiffres
0 et 1 utilisés pour représenter l'information s'appellent des {\em bits}. Nous allons
commencer par étudier d'au\-tres formes de codages plus anciennes.

\section{Code Morse}
Le télégraphe électrique, ancêtre de notre téléphone, fut inventé en 1837 par 
l'américain \href{https://fr.wikipedia.org/wiki/Samuel_Morse}{Samuel Morse} \cite{bertho}. Le télégraphiste émetteur dispose d'un manipulateur qui est un simple levier
à ressort permettant d'envoyer des signaux courts (appelés points) et des signaux longs (appelés
traits) séparés par des silences ou blancs (correspondant à l'absence de signal). Le signal est 
transmis le long d'un fil électrique, le retour se faisant par la terre. Chaque
mouvement  du manipulateur ferme le circuit et provoque un mouvement bref ou long du récepteur 
qui reproduit les signaux
sur une bande de papier.  Pour envoyer un message, on commence par l'abréger en {\em style
télégraphique} puis on traduit chaque signe (lettre, chiffre, signe de ponctuation, \ldots) en
une suite de points et de traits séparés par des blancs selon le {\em code Morse} de la figure 
\ref{fig:codemorse} (page \pageref{fig:codemorse}). On laisse sept blancs à la fin de chaque mot. 
Divers autres codes sont utilisés pour indiquer le début et la fin du message, signaler les erreurs,
etc. Le télégraphe de Morse permettait de transmettre environ 25 mots à la minute.
\begin{figure}[t]
\vspace{9.3cm}%
\includePostScript{CodeMorse}{0.95}{0.95}%
% \vspace{8.8cm} %8.9cm donne la figure sur une seule page !
 \caption{Code Morse}\label{fig:codemorse}
\end{figure}

\vspace{5.0cm}

\section{Numération}
\index{Numération décimale}
Notre codage des entiers naturels (positifs ou nuls) est une numération positionnelle et décimale.
C'est une numération décimale car nous utilisons les dix chiffres 0, 1, \ldots, 9 d'origine
indo-arabe. C'est une numération positionnelle car chaque chiffre dans un nombre représente une
valeur qui dépend de sa position dans le nombre. Par exemple dans 1789, le 1 représente 1000,
le 7 représente 700, le 8 représente 80 et le 9 représente 9. On a donc $1789 = 1000 + 700
+ 80 + 9 = (1 \times 10^{3}) +  (7 \times 10^{2}) +  (8 \times 10^{1}) + (9 \times 10^{0})$.  
Pour les chiffres après la virgule on considère les inverses. Par exemple $0.147= (1\times 10^{-1})+
(4\times 10^{-2})+(7\times 10^{-3})=\frac{1}{10}+\frac{4}{100}+\frac{7}{1000}=\frac{100}{1000}+\frac{40}{1000}+\frac{7}{1000}=\frac{147}{1000}$.
Les
mathématiciens anciens et modernes ont inventé bien d'autres façons de coder les nombres. 
La figure \ref{fig:code1-9} (page \pageref{fig:code1-9}) montre différentes façons de coder le nombre
123 456 789, certaines remontant aux civilisations antiques \cite{guitel}. La supériorité du
codage positionnel et décimal des entiers naturels ne tient pas au choix de la base dix mais
à la facilité avec laquelle on peut faire les opérations
arithmétiques usuelles, par exemple en utilisant un boulier (longtemps en usage et maintenant disparu des écoles
communales)

\begin{figure}[htbp]
 \includevspacePostScript[18.0cm]{Codages123456789}{0.95}{0.95}
 \caption{Codages de 123 456 789 en numérations babylonienne, égyptien-\label{fig:code1-9}}
 \centering{ne, grecque, romaine, hexadécimale, décimale, octale, ternaire et binaire}
\end{figure}


\section{Codage binaire}

Toutes les informations rangées dans la mémoire de l'ordinateur ou sur disquette doivent
être représentées par des suites de zéros et de uns. Sur les micro-ordinateurs,
tels qu'ils sont conçus actuellement, la longueur de ces suites de bits doit être multiple
de 8. Selon leur longueur on leur donne des noms spéciaux comme {\em octet} pour une suite
de 8 bits, {\em demi-mot} pour une suite de 16 bits, {\em mot} pour 32 bits et {\em double-mot}
pour 64.

Pour montrer que toute information peut être représentée par une suite de bits (à condition
de connaître le code) nous allons prendre l'exemple des entiers puis celui des caractères et 
des textes tapés au clavier.

\subsection{Entiers naturels}

\index{Numération binaire}
Pour représenter les entiers positifs par une suite de bits, il suffit de les coder en base
2. Par exemple~:
\[\begin{array}{lll}
147 & = &128 +  16 + 2 + 1 \\
       & = & (1 \times 2^{7}) + (0 \times 2^{6}) + (0 \times 2^{5}) + (1 \times 2^{4}) + (0 \times 2^{3}) + (0 \times 2^{2}) \\
       &    & + (1 \times 2^{1}) + (1 \times 2^{0})
\end{array}\]
s'écrit~: \[10010011\] en base deux. De manière plus générale, le
 nombre qui s'écrit~: \[a_{n}a_{n-1} \ldots a_{1}a_{0}\]
 en base deux (où $a_{n} \ldots a_{0}$ valent 0 ou 1)
 vaut~:
\[(a_{n} \times 2^{n}) + (a_{n-1} \times 2^{n-1}) + \ldots 
+ (a_{1} \times 2^{1}) + (a_{0} \times 2^{0})\]
(où $2^{n}$ est égal à $2$ multiplié $n$ fois par lui-même si $n>0$ et $2^{0} = 1$).
Avec cette convention on peut donc coder
 sur un octet les entiers naturels compris entre 0 (00000000) et 255 (11111111).

\subsection{Caractères (code {\sc ascii})}
\index{Caractère}
 Les programmes \pascal\ que nous écrivons sont des suites de caractères. À chaque fois
 que nous tapons sur une touche du clavier le code correspondant du caractère est entré dans la
 mémoire de l'ordinateur. Les caractères cor\-res\-pon\-dant aux chiffres, aux lettres minuscules 
 ou majuscules et aux caractères usuels comme $+$, $\ast$,
 $($, $)$, \ldots  s'appellent
 les caractères ``imprimables''. Certains autres caractères qui
ne peuvent pas s'écrire s'appellent les caractères ``spéciaux.'' C'est le cas par
exemple des ``ca\-rac\-tè\-re de contrôle'' et ``ca\-rac\-tè\-re d'échappement''
que l'on obtient respectivement en tapant les touches
%touche Ctrl
\begin{picture}(22,12)(0, 2)
\put(0,0){\framebox(22,12)[lb]{}}
\put(1.5,2){\makebox(0,0)[lb]{\tiny Ctrl}}
\end{picture}\hspace{0.3ex}
et
%Esc (touche d'echappement)
\begin{picture}(14,12)(0, 2) 
\put(0,0){\framebox(14,12)[lb]{}}
\put(1.5,2){\makebox(0,0)[lb]{\tiny Esc}}
\end{picture}\hspace{0.3ex}. 
Ces caractères (imprimables ou spéciaux) sont souvent représentés sur un octet (8 bits) selon
un code international qui s'appelle le code {\sc ascii}.  Les caractères imprimables dont
le code {\sc ascii} est compris entre 32 et
126 sont donnés dans la table \ref{tab:ascii} ci-dessous. Par exemple, le code {\sc ascii} 65 du
caractère {\cha $\scriptstyle '$A$\scriptstyle '$} s'obtient dans cette table en ajoutant le numéro 64 de la ligne
où se trouve le caractère {\cha $\scriptstyle '$A$\scriptstyle '$} au numéro 1 de la colonne
de ce caractère.

\index{Code {\sc ascii}}\begin{table}[htbp]
\begin{center}
\begin{tabular}{@{}|c|c|c|c|c|c|c|c|c|c|c|c|c|c|c|c|c|@{}} \cline{2-17}
\multicolumn{1}{c|}{$+$}       &  \tiny \ 0  & \tiny \ 1 & \tiny \ 2 &\tiny \ 3 &\tiny \ 4 &\tiny \ 5 &\tiny \ 6 &\tiny \ 7
        &\tiny \ 8 &\tiny \ 9 & \tiny 10 &\tiny 11&\tiny 12 &\tiny 13 &\tiny 14 &\tiny 15 \\ \hline
\tiny 32    &  \blanc & \texttt{!} & \texttt{"} & \texttt{\#} & \texttt{\$} & \texttt{\%} & \texttt{\&} & \texttt{'} 
        & \texttt{(} & \texttt{)} & \texttt{*} & \texttt{+} & \texttt{,} & \texttt{-} & \texttt{.} & \texttt{/} \\ \hline
\tiny 48    & \texttt{0} & \texttt{1} & \texttt{2} & \texttt{3} & \texttt{4} & \texttt{5} & \texttt{6} & \texttt{7} 
        & \texttt{8} & \texttt{9} & \texttt{:} & \texttt{;} & \texttt{<} & \texttt{=} & \texttt{>} & \texttt{?} \\ \hline
\tiny 64    & \texttt{@} & \texttt{A} & \texttt{B} & \texttt{C} & \texttt{D} & \texttt{E} & \texttt{F} & \texttt{G} 
         & \texttt{H} & \texttt{I} & \texttt{J} & \texttt{K} & \texttt{L} & \texttt{M} & \texttt{N} & \texttt{O} \\ \hline
\tiny 80     & \texttt{P} & \texttt{Q} & \texttt{R} & \texttt{S} & \texttt{T} & \texttt{U} & \texttt{V} & \texttt{W} 
         & \texttt{X} & \texttt{Y} & \texttt{Z} & \texttt{[} & \texttt{$\backslash$} & \texttt{]} & \texttt{\^{ } } & \texttt{\_} \\ \hline
\tiny 96     & \texttt{`} & \texttt{a} & \texttt{b} & \texttt{c} & \texttt{d} & \texttt{e} & \texttt{f} & \texttt{g} 
         & \texttt{h} & \texttt{i} & \texttt{j} & \texttt{k} & \texttt{l} & \texttt{m} & \texttt{n} & \texttt{o} \\ \hline
\tiny 112  & \texttt{p} & \texttt{q} & \texttt{r} & \texttt{s} & \texttt{t} & \texttt{u} & \texttt{v} & \texttt{w} 
         & \texttt{x} & \texttt{y} & \texttt{z} & \texttt{\{} & \texttt{|} & \texttt{\}} & \texttt{\~{ }} &  \multicolumn{1}{c}{}\\ \cline{1-16}
\end{tabular}
\end{center}
\caption{Caractères de codes {\sc ascii} compris entre 32 et 126\label{tab:ascii}}
\end{table}

Quand nous disons que le caractère {\cha $\scriptstyle '$A$\scriptstyle '$} a pour code {\sc ascii} 65, il faut comprendre 
que ce caractère est représenté dans la mémoire de l'ordinateur par l'octet
01000001 qui représente également le nombre 65 écrit en base deux. On remarquera que 
le caractère zéro a pour code 48. L'entier 0 (codé en machine sur un octet par 00000000) a donc
un code dif\-fé\-rent du ca\-rac\-tè\-re zé\-ro {\cha $\scriptstyle '$0$\scriptstyle '$} (codé en machine par 00110000). La table
\ref{tab:asciispec} (page \pageref{tab:asciispec}) montre les codes {\sc ascii} de quelques 
caractères spéciaux.

 \begin{table}[htbp]
\begin{center}
\begin{tabular}{@{}|c|c|c|c|c|@{}} \hline
Code {\sc ascii} & 8 &  13 & 27 & 127 \\ \hline
Touche du clavier \rule[-3mm]{0mm}{8mm}& 
%touche d'effacement 8
\framebox[11.mm][l]{$\leftarrow$~~~}
&
\framebox[11.mm][l]{$\hookleftarrow$~~~}
&  
%Esc (touche d'echappement) 27
\begin{picture}(14,12)(0, 2) 
\put(0,-1){\framebox(14,12)[lb]{}}
\put(1.5,2){\makebox(0,0)[lb]{\tiny Esc}}
\end{picture}
&  
%touche Del
\begin{picture}(14,12)(0, 2)
\put(0,-1){\framebox(14,12)[lb]{}}
\put(1.5,2){\makebox(0,0)[lb]{\tiny Del}}
\end{picture}
\\ \hline
\end{tabular}
\end{center}
\caption{Code {\sc ascii} de quelques caractères spéciaux\label{tab:asciispec}}
\end{table}

\subsection{Texte}

Un texte (comme par exemple un programme \pascal) est une suite de ca\-rac\-tè\-res. Le code
d'un texte en machine est la suite des codes des caractères qui le composent. Par exemple
le texte {\cha $\scriptstyle '$Pascal$\scriptstyle '$} sera représenté en machine par la suite des codes {\sc ascii} des caractères
 {\cha $\scriptstyle '$P$\scriptstyle '$}, {\cha $\scriptstyle '$a$\scriptstyle '$}, {\cha $\scriptstyle '$s$\scriptstyle '$},
 {\cha $\scriptstyle '$c$\scriptstyle '$}, {\cha $\scriptstyle '$a$\scriptstyle '$}, {\cha $\scriptstyle '$l$\scriptstyle '$}
  soit, d'après la table  \ref{tab:ascii}
(page \pageref{tab:ascii}), 
la suite 80 (01010000), 
97 (01100001), 115 (01110011), 99 (01100011), 97 (01100001), 108 (01101100).
Le code du texte {\cha $\scriptstyle '$Pascal$\scriptstyle '$} est donc
01010000 01100001 01110011 01100011 01100001 01101100.

\subsection{Chaînes de caractères}
\index{Chaîne de caractères}

Un programme permet de faire manipuler par l'ordinateur des booléens, des entiers mais
également des textes, tous codés en binaire. Pour ne pas confondre le texte du programme
et les textes manipulés par le programme, nous appelerons ces derniers des ``chaînes de 
caractères''.

\subsubsection{Usage et notation en \protect\pascal~:}
En \pascal\ on peut manipuler des textes limités à 255 caractères que l'on appelle des
``chaînes de caractères'' ({\idt string} (prononcer stri-ngue) en anglais). Un paramètre
formel de procédure {\idt X} dont la valeur est une chaîne de ca\-rac\-tè\-res doit être
déclaré {\idt X $:$ string}. Le paramètre effectif correspondant est une chaîne de
caractères, comme par exemple~:
\begin{center}
{\cha $\scriptstyle '$Ceci\ est\ une\ chaine\ de\ caracteres\idt $\scriptstyle '$}
\end{center}
qui s'écrit entre apostrophes. S'il est facile de distinguer un identificateur ou un entier dans le texte d'un
programme, il est plus difficile de reconnaître un texte dans ce texte$\:$! L'usage des apostrophes permet donc de 
délimiter la partie du texte constituant la chaîne de caractères de celle du reste du programme. Mais ceci pose
problème si la chaîne de caractères contient une apostrophe. C'est pourquoi on convient de doubler les apostrophes
dans une chaîne de caractères écrite entre apostrophes, 
comme dans~:
\begin{center}
{\cha $\scriptstyle '$Le\ langage\ \idt $\scriptstyle '$\cha $\scriptstyle '$pascal\idt $\scriptstyle '$\cha $\scriptstyle '$\idt $\scriptstyle '$}.
\end{center}
Une chaîne de caractères réduite à une seule apostrophe s'écrit donc
{\cha $\scriptstyle '$\idt $\scriptstyle '$\cha $\scriptstyle '$\idt $\scriptstyle '$}.

\subsubsection{Concaténation de chaînes de caractères~:}
\index{Concaténation}
Si $c_{1}$, $c_{2}$, \ldots, $c_{n}$ sont des chaînes de
caractères alors {\idt concat$(c_{1}$, $c_{2}$, \ldots, $c_{n})$} est la chaîne de
caractères obtenue en concaténant (c'est-à-dire en mettant bout à bout) les chaînes de
caractères $c_{1}$, $c_{2}$, \ldots, $c_{n}$. Par exemple~:
\begin{center}
{\idt concat$($\cha $\scriptstyle '$Le\ \idt $\scriptstyle ',$\ \cha $\scriptstyle '$langage\ \idt $\scriptstyle ',$\ %
\cha $\scriptstyle '$\idt $\scriptstyle '$\cha $\scriptstyle '$\idt $\scriptstyle ',$\ %
\cha $\scriptstyle '$pascal\idt $\scriptstyle ',$\ \cha $\scriptstyle '$\idt $\scriptstyle '$%
\cha $\scriptstyle '$\idt $\scriptstyle ')$}
\end{center}
est égal à la chaîne de caractères qui s'écrit également~:
\begin{center}
{\cha $\scriptstyle '$Le\ langage\ \idt $\scriptstyle '$\cha $\scriptstyle '$pascal\idt $\scriptstyle '$\cha $\scriptstyle '$\idt $\scriptstyle '$}
\end{center}

\subsubsection{Écriture sur le dessin et sous le dessin du robot~:}
Que son crayon soit levé ou baissé, le \robot\ peut écrire des chaînes de ca\-rac\-tè\-res,  
sans se déplacer et sans changer son orientation. Il peut écrire dans la fenêtre de dessin ou dans le cadre d'affichage des
messages situé sous la fenêtre de dessin.

La commande {\idt Ecrire$(c)$} indique au \robot\ qu'il faut écrire la chaîne de caractères $c$ à droite du point
où il se trouve dans sa fenêtre de dessin.
La commande {\idt Message$(c)$} indique au \robot\ qu'il faut écrire la chaîne de ca\-rac\-tè\-res $c$ dans le
cadre prévu en dessous de la fenêtre de déplacement. Ceci a pour effet d'effacer le message précédent 
affiché dans ce
cadre. Pour laisser à l'utilisateur le temps de lire le message, on peut, si nécessaire, utiliser la procédure
{\idt  MarquerUnePause} qui attend pour continuer que l'utilisateur tape sur une touche ou clique sur le bouton de la 
souris. La commande {\idt  EffacerMessage} efface ce message. Le choix d'une police, d'un style ou d'une taille de
caractères autre que le choix standard est expliqué sur les disquettes d'accompagnement du livre.

Si $n$ est un entier alors {\idt EntierEnChaine$(n)$} est la chaîne de caractères qui dénote cet entier dans la notation
décimale usuelle. Si $b$ est un booléen alors  {\idt BooleenEnChaine$(b)$} est 
{\cha $\scriptstyle '$Vrai\idt $\scriptstyle '$}
si $b$ est égal à {\idt true} et
{\cha $\scriptstyle '$Faux\idt $\scriptstyle '$}
si $b$ est égal à {\idt false}.  
Par exemple si {\idt X} est un paramètre entier de procédure égal à 17 alors la
commande~:
\begin{center}
{\idt  ecrire$($concat$($\cha $\scriptstyle '$La\ valeur\ de\ X\ est\ \idt $\scriptstyle ',$\ EntierEnChaine$($X$),$\ \cha $\scriptstyle '$.\idt $\scriptstyle '))$}
\end{center}
\' ecrit la concaténation des chaînes  de caractères
{\cha $\scriptstyle '$La\ valeur\ de\ X\ est\ $\scriptstyle '$},
{\cha $\scriptstyle '$17$\scriptstyle '$}
et {\cha $\scriptstyle '$.$\scriptstyle '$}, c'est-à-dire~:
\begin{center}
{\idt La valeur de X est 17.}
\end{center}
à la droite du \robot. 
Si {\idt B} est un paramètre booléen de procédure égal à {\idt true} alors la
commande~:
\begin{center}
{\idt  message$($concat$($\cha $\scriptstyle '$B\ =\ \idt $\scriptstyle ',$\ BooleenEnChaine$($B$)));$}
\end{center}
écrit la concaténation des chaînes  de caractères
{\cha $\scriptstyle '$B\ =\ $\scriptstyle '$}
et {\cha $\scriptstyle '$Vrai$\scriptstyle '$}, c'est-à-dire~:
\begin{center}
{\idt B = Vrai}
\end{center}
dans le cadre des messages.



\section{Musique}

Ce paragraphe est destiné aux musiciens qui connaissent le solfège, c'est-à-dire la manière 
de coder la musique.  Le \robot\ peut jouer de la musique, de 
manière il est vrai assez limitée, puisqu'il ne peut jouer du piano qu'à un seul doigt. Pour
coder la musique pour le \robot\ nous disposons des procédures avec paramètres entiers. Il
nous suffit donc d'établir la correspondance avec la théorie classique de la musique
\cite{solfege} pour définir comment écrire de la musique pour le \robot.

\subsection{Figures de notes}

La procédure {\idt Note$($h$,$ d$)$} permet
de faire entendre sur le haut-parleur du micro-ordinateur une note de hauteur {\idt h} et de
figure {\idt d}. La hauteur de la note est codée par un entier {\idt h} comme indiqué à la figure
\ref{fig:codenotes} (page \pageref{fig:codenotes}).
\begin{figure}[htbp]
%\includevspacePostScript[13.4cm]{CodageNotes}{0.95}{0.95}
 \begin{flushleft}
\parindent0pt
\noindent Octave 2 :\\
\begin{music}
\parindent0pt
\startextract
\NOtes\ql{q}\ql{^q}\ql{r}\ql{^r}\ql{s}\ql{t}\ql{^t}\ql{u}\ql{^u}\ql{v}\ql{^v}\ql{w}\ql{x}\enotes
\endextract
\end{music}
\vspace*{-1.21em}%
\def\sauth{\hspace*{1.19em}}%
\hspace*{17mm}{\small24\sauth25\sauth26\sauth27\sauth28\sauth29\sauth30\sauth31\sauth32\sauth33\sauth34\sauth35\sauth36}\\[1ex]
\noindent Octave 1 :\\
\begin{music}
\parindent0pt
\startextract
\NOtes\ql{j}\ql{^j}\ql{k}\ql{^k}\ql{l}\ql{m}\ql{^m}\ql{^n}\ql{n}\ql{o}\ql{^o}\ql{p}\ql{q}\enotes
\endextract
\end{music}
\vspace*{-1.em}%
\def\sauth{\hspace*{1.19em}}%
\hspace*{17mm}{\small12\sauth13\sauth14\sauth15\sauth16\sauth17\sauth18\sauth19\sauth20\sauth21\sauth22\sauth23\sauth24}\\[1ex]
\noindent Octave 0 :\\
\begin{music}
\parindent0pt
\startextract
\NOtes\qu{c}\qu{^c}\qu{d}\qu{^d}\qu{e}\qu{f}\qu{^f}\qu{g}\qu{^g}\qu{h}\qu{^h}\ql{i}\ql{j}\enotes
\endextract
\end{music}
\vspace*{-0.75em}%
\def\sauth{\hspace*{1.67em}}%
\hspace*{18mm}{\small0\sauth1\sauth2\sauth3\sauth4\sauth5\sauth 6\sauth7\sauth8\sauth9\sauth1\rlap{0}\sauth1\rlap{1}\sauth1\rlap{2}}\\[1ex]
\noindent Octave -1 :\\
\begin{music}
\parindent0pt
\setclef1\bass
\startextract
\NOtes\qu{J}\qu{^J}\qu{K}\qu{^K}\qu{L}\qu{M}\qu{^M}\qu{N}\qu{^N}\qu{a}\qu{^a}\ql{b}\ql{c}\enotes
\endextract
\end{music}
\vspace*{-1.2em}%
\def\sauth{\hspace*{1.445em}}%
\hspace*{16mm}{\small-1\rlap{2}\sauth-1\rlap{1}\sauth-1\rlap{0}\sauth -9\sauth -8\sauth -7\sauth 6\sauth -5\sauth -4\sauth -3\sauth -2\sauth -1\sauth 0}\\[1ex]
\noindent Octave -2 :\\
\begin{music}
\parindent0pt
\setclef1\bass
\startextract
\NOtes\qu{C}\qu{^C}\qu{D}\qu{^D}\qu{E}\qu{F}\qu{^F}\qu{G}\qu{^G}\qu{H}\qu{^H}\qu{I}\qu{J}\enotes
\endextract
\end{music}
\def\lignes#1#2{\begin{tabular}[t]{@{}c@{}}#1\\#2\end{tabular}}
\def\sauth{\hspace*{0.86em}}%
\hspace*{16mm}{\small\lignes{-24}{do}\sauth -23\sauth \lignes{-22}{r\'e}\sauth-21\sauth \lignes{-20}{mi}\sauth\lignes{-18}{fa}\sauth -18\sauth \lignes{-17}{sol}\sauth -16\sauth\lignes{-15}{la}\sauth -14\sauth\lignes{-13}{si}\sauth\lignes{-12}{do}}
\end{flushleft}%
 \caption{Codage de la hauteur des notes de musique par des entiers}\label{fig:codenotes}
\end{figure}
La figure de la note indique sa durée qui est codée par une constante {\idt d} qui peut être {\idt ronde}, 
{\idt ronde\_pointee}, 
{\idt blanche}, 
{\idt blanche\_pointee}, 
{\idt blanche\_en\_triolet},  
{\idt noire}, 
{\idt noi\-re\_poin\-tee},       
{\idt noi\-re\_\-en\_\-tri\-o\-let},      
 {\idt croche}, 
 {\idt croche\_pointee}, 
{\idt cro\-che\_\-en\_\-trio\-let},  
{\idt double\_croche},       
{\idt dou\-ble\_\-cro\-che\_\-poin\-tee}, 
{\idt double\_croche\_en\_trio\-let},    
{\idt tri\-ple\_\-cro\-che},   
{\idt tri\-ple\_\-cro\-che\_\-poin\-tee}, 
{\idt tri\-ple\_\-cro\-che\_\-en\_\-tri\-o\-let},   {\idt qua\-dru\-ple\_\-cro\-che}           ou 
{\idt quadruple\_croche\_en\_triolet}. La durée relative d'une {\idt ron\-de\_\-poin\-tee} est 
une fois et demie celle d'une
 {\idt ronde}. La durée relative d'une {\idt ronde} est  deux fois celle d'une
 {\idt blanche}, trois fois celle d'une {\idt blan\-che\_\-en\_\-trio\-let}, \ldots   
 
 \subsection{Mouvement}
 
 Le mouvement, 
 c'est-à-dire le degré de vitesse ou de lenteur avec lequel doit être exécuté le morceau de 
 musique détermine la durée absolue des différentes figures de notes. Ce mouvement est défini
 par la procédure {\idt metronome$(n)$} qui indique le nombre $n$ d'oscillations du
 métronome par minute, sachant qu'une noire est exécutée à chaque battement.  Par exemple
 {\idt metronome$($60$)$} indique qu'il faut jouer 60 noires à la minute. Les mouvements
 sont gé\-né\-ra\-le\-ment définis par des termes italiens (adagio, allegro, prestissimo,\ldots) que
 l'on place au début d'un morceau et au dessus de la portée. On peut les définir
 en \pascal\ à l'aide des déclarations de constantes suivantes~:

{
\parindent \parindentprog
\medskip

\pagebreak[3] 
\vspace{\neglskip}\mbox{\blc  \idt  \mcl const\idt }\par 
\vspace{\neglskip}\mbox{\blc  \blk\blk\idt  largo\ $=$\ 50$;$}\par 
\vspace{\neglskip}\mbox{\blc  \blk\blk\idt  larghetto\ $=$\ 63$;$}\par 
\vspace{\neglskip}\mbox{\blc  \blk\blk\idt  adagio\ $=$\ 71$;$}\par 
\vspace{\neglskip}\mbox{\blc  \blk\blk\idt  andante\ $=$\ 92$;$}\par 
\vspace{\neglskip}\mbox{\blc  \blk\blk\idt  moderato\ $=$\ 114$;$}\par 
\vspace{\neglskip}\mbox{\blc  \blk\blk\idt  allegro\ $=$\ 144$;$}\par 
\vspace{\neglskip}\mbox{\blc  \blk\blk\idt  presto\ $=$\ 184$;$}\par 
\vspace{\neglskip}\mbox{\blc  \blk\blk\idt  prestissimo\ $=$\ 204$;$}\par 

\medskip }

\subsection{Notes}

\noindent Le numérotage des notes de la figure \ref{fig:codenotes} (page \pageref{fig:codenotes}) n'est
pas facile à mé\-mo\-ri\-ser. C'est pourquoi nous allons coder les notes par 
une expression entière portant sur des constantes dont il est facile de se souvenir.

Une note est définie par son nom $n$ (à savoir do, ré, mi, fa, sol, la, si), son éventuelle
altération $a$ (c'est-à-dire un signe qui modifie le son de la note à laquelle il est affecté
qui peut être le bémol $\flat$ (qui abaisse le son de la note d'un demi-ton),  
le dièse $\sharp$ (qui élève le son de la note d'un demi-ton) ou
le bécarre $\natural$ (qui détruit l'effet du dièse ou du bémol)) et
son octave $o$ (qui indique à quelle série de notes portant le même nom appartient cette note). 
Nous codons les octaves du grave à l'aigu par des entiers -3, -2, \ldots, 3, 4 , 
l'octave 0 commençant par le do au centre du clavier du piano. Définissons les
constantes suivantes en \pascal\footnote{\ On remarquera que la note {\mcl do} correspond au mot 
anglais utilisé dans la boucle ``for''.
Ce mot ne pouvant pas être utilisé pour un autre usage en \pascal\ nous sommes dans l'obligation
de trouver un autre nom pour la note do. Nous laissons le choix entre {\idt ut} ou {\idt d0} (lettre d
suivie du chiffre zéro).}~:

{
\parindent \parindentprog
\medskip

\pagebreak[3] 
\vspace{\neglskip}\mbox{\blc  \idt  \mcl const\idt }\par 
\vspace{\neglskip}\mbox{\blc  \blk\idt  \cmt  \{\ notes\ :\ \}\idt }\par 
\vspace{\neglskip}\mbox{\blc  \blk\blk\idt  ut\ $=$\ 0$;$}\par 
\vspace{\neglskip}\mbox{\blc  \blk\blk\idt  d0\ $=$\ 0$;$}\par
\vspace{\neglskip}\mbox{\blc  \blk\blk\idt  re\ $=$\ 2$;$}\par 
\vspace{\neglskip}\mbox{\blc  \blk\blk\idt  mi\ $=$\ 4$;$}\par 
\vspace{\neglskip}\mbox{\blc  \blk\blk\idt  fa\ $=$\ 5$;$}\par 
\vspace{\neglskip}\mbox{\blc  \blk\blk\idt  sol\ $=$\ 7$;$}\par 
\vspace{\neglskip}\mbox{\blc  \blk\blk\idt  la\ $=$\ 9$;$}\par 
\vspace{\neglskip}\mbox{\blc  \blk\blk\idt  si\ $=$\ 11$;$}\par 
\vspace{\neglskip}\mbox{\blc  \blk\blk\idt  \cmt  \{\ altérations\ :\ \}\idt }\par 
\vspace{\neglskip}\mbox{\blc  \blk\blk\idt  bemol\ $=$\ $-$1$;$}\par 
\vspace{\neglskip}\mbox{\blc  \blk\blk\idt  becarre\ $=$\ 0$;$}\par 
\vspace{\neglskip}\mbox{\blc  \blk\blk\idt  diese\ $=$\ 1$;$}\par 
\vspace{\neglskip}\mbox{\blc  \blk\blk\idt  \cmt  \{\ octaves\ :\ \}\idt }\par 
\vspace{\neglskip}\mbox{\blc  \blk\blk\idt  octave\ $=$\ 12$;$}\par 

\medskip }

\noindent Le numéro de la note tel qu'il est  défini à la figure \ref{fig:codenotes} 
(page \pageref{fig:codenotes}) est donné par la valeur de l'expression entière 
$(((o \ast {\idt octave}) + n) + a)$. Si la note n'est pas altérée, son code est simplement
$((o \ast {\idt octave}) + n)$. Par exemple le do du milieu du clavier du piano correspond à 
$o = 0$, $n = ${\idt d0} soit $((0 \ast ${\idt octave})$ + 0) = 0$. Le ré dièse, une octave plus haut
a pour code $(((1 \ast ${\idt octave}$) + ${\idt re}$) + ${\idt diese}$)
= (((1 \ast 12) + 2) + 1) = 15$. Le ré bémol voisin
a pour code $(((1 \ast ${\idt octave}$) + ${\idt re}$) + ${\idt bemol}$)
= (((1 \ast 12) + 2) + (-1)) = 13$.

\subsection{Silences}

La procédure {\idt silence$(s)$} permet d'interrompre le son pour une durée re\-la\-ti\-ve fixée par
la figure de silence $s$ qui peut être l'une des constantes  
  {\idt pause},
  {\idt pause\_pointee},
  {\idt demi\_pause\_pointee},
  {\idt demi\_pause},
  {\idt soupir},
  {\idt sou\-pir\_\-poin\-te}\footnote{\ Il n'est pas d'usage de pointer le soupir.},
  {\idt de\-mi\_\-sou\-pir},
  {\idt demi\_soupir\_pointe},
  {\idt quart\_de\_soupir},
  {\idt quart\_de\_soupir\_pointe},
  {\idt huitieme\_de\_soupir},
  {\idt huit\-ieme\_\-de\_\-sou\-pir\_\-poin\-te}
ou  {\idt seizieme\_de\_soupir}.

\subsection{Nuances}

Les nuances sont les différents degrés de force par lesquels peuvent passer un ou plusieurs sons, un trait
ou un morceau entier. La procédure {\idt nuance$(n)$} définit la nuance avec laquelle doit être jouée
la suite du morceau. Les termes $n$ de nuances possibles sont, du plus faible au plus fort, les suivants~:
{\idt pia\-nis\-si\-mo},
{\idt piano},
{\idt mezzo\_piano},
{\idt un\_poco\_piano},
{\idt sotto\_voce},
{\idt mezza\_voce},
{\idt un\_\-po\-co\_\-for\-te},
{\idt mezzo\_forte},
{\idt forte} ou
{\idt fortissimo}.

\begin{exemple}[Au clair de la lune]{\rm
``Au clair de la lune'', qui aurait été composée par le musicien Jean-Baptiste Lulli,
est sans doute la plus célèbre des chansons enfantines françaises~:

%\includevspacePostScript[4.7cm]{AuClairDeLaLune}{0.95}{0.95}
%%&frLaTeX
%\documentclass{article}
%\usepackage{musixtex}
%\begin{document}
\begin{flushleft}
\parindent0pt
\begin{music}
\generalmeter{\meterC}
\generalsignature{1}
\nostartrule
\parindent0pt%
\nobarnumbers
\startpiece
\notesp\qu{g}\qu{g}\qu{g}\qu{h}\enotes\bar
\notesp\hl{i}\hu{h}\enotes\bar
\notesp\qu{g}\ql{i}\qu{h}\qu{h}\enotes\Setvolta1\bar
\notesp\zw{g}\enotes\setendvoltabox\Setvolta2\rightrepeat%
\notesp\zw{g}\enotes\setendvoltabox\alaligne
\NOtesp\qu{h}\qu{h}\qu{h}\qu{h}\enotes\bar
\NOtesp\hu{e}\hu{e}\enotes\bar
\NOtesp\qu{h}\ql{g}\qu{f}\qu{e}\enotes\bar
\NOtesp\zw{d}\enotes\alaligne
\NOtesp\qu{g}\qu{g}\qu{g}\qu{h}\enotes\bar
\NOtesp\hl{i}\hu{h}\enotes\bar
\NOtesp\qu{g}\ql{i}\qu{h}\qu{h}\enotes\bar
\NOtesp\zw{g}\enotes
\Endpiece
\end{music}
\end{flushleft}%
%\end{document}%

\noindent Le programme ci-dessous permet de la faire jouer par l'ordinateur~:

 \input AuClairDeLaLune.tex
\noindent $\Box$ }\end{exemple}

\noindent \rule{137mm}{0.1mm} %---------------------------------------------
\begin{exercices}{\ \\
\vspace{-\baselineskip}

\tiret  Écrire des procédures {\idt DebutTransmission}, {\idt A}, \ldots, {\idt Z}, {\idt 0},
\ldots, {\idt 9}, {\idt Po} (point), {\idt Vg} (virgule),
 {\idt Pv} (point-virgule), {\idt Dp} (deux-points), {\idt Pi} (point d'in\-ter\-ro\-ga\-tion),
 {\idt Pe} (point d'exclamation), {\idt Tu} (trait d'union), {\idt Ap} (apostrophe),
  {\idt Di} (barre de division),
 {\idt Erreur}, {\idt FinTransmission}, {\idt Fm} (fin de mot) et
 {\idt ALaLigne} (passer à la ligne suivante) puis les utiliser pour coder
 des messages en morse, chaque signe du message en clair étant codé en appelant l'une
 de ces procédures. Par exemple ``{\em message codé}'' sera traduit par la suite d'instructions
`` {\idt M$;$\ E$;$\ S$;$\ S$;$\ A$;$\ G$;$\ E$;$\ Fm$;$\ ALaLigne
 \ C$;$\ O$;$\ D$;$\ E$;$\ Fm$;$}''. On pourra ajouter le son en utilisant les procédures suivantes
 du \robot~:
 \begin{list}{$\scriptscriptstyle\bullet$}{
}
   \item {\idt Bip} produit un son court$\:$;
   \item {\idt Biiip} produit un son long$\:$;
   \item {\idt Delai$(${\rm\em t}$)$} fait attendre pendant {\em t} secondes.
 \end{list}
 
 \tiret Avant l'invention de l'électricité, on utilisait en France le télégraphe  aérien inventé en
 1793 par le citoyen Claude Chappe. Le télégraphe est constitué par trois bras de bois articulés
 à l'extrémité d'un poteau ins\-tal\-lé sur une tour au sommet d'une colline ou sur un clocher d'église.
 Dès qu'il aperçoit à la lunette un signe d'activité au poste le plus proche, le {\em télégraphier}
 s'empresse de manœuvrer à son tour les bras de sa machine pour transmettre (sans le comprendre)
  le message au poste suivant. De proche en proche le message arrive à destination. Il s'agit en
  quelque sorte d'une version perfectionnée des signaux de fumée des Indiens$\:$! Après une première
  expérience le 12 juillet 1793, la première ligne fut ouverte entre Paris et Lille le 17 août 1794
  \cite{bertho}. Les signaux du télégraphe optique de Chappe pouvaient être décodés
  selon plusieurs codes. Le code alphabétique est donné à la figure \ref{fig:chappe} (page
  \pageref{fig:chappe}). Écrire un programme qui permet d'afficher successivement, un par un, les
  signaux et envoyer le message ``Paris est tranquille et les bons citoyens sont contents''
  que Bonaparte fit diffuser immédiatement après le coup d'État du 18 brumaire (9-10 novembre 
  1799).
  
  Pour laisser apparaître le signal un certain temps on utilisera la pro\-cé\-du\-re 
  {\idt Delai$(${\rm\em t}$)$} qui fait attendre pendant {\rm\em t} secondes. Pour effacer l'écran
  entre chaque signal, on utilisera la procédure {\idt cf$($blanc$)$} qui peint le fond de l'écran
  en blanc (ou dans une autre couleur selon le paramètre) et place le \robot\ dans les conditions initiales
  (visible au centre de l'écran, or\-ien\-té vers le nord, crayon noir d'épaisseur 1 baissé, avec une grille 
  de déplacement (12~$\times$~12)).

  \tiret Écrire un programme Pascal de traduction des lettres de l'alphabet en alphabet phonétique
  selon le code international de la table \ref{tab:alphone} (page \pageref{tab:alphone}).

   \tiret Écrire un programme Pascal de cryptage et de décryptage de messages au moyen d'un code secret
   simple, obtenu par exemple par permutation des lettres de l'alphabet.

\noindent $\Box$}\end{exercices}
\begin{exercices}{\ \\
\vspace{-\baselineskip}

\tiret Quelle est la représentation binaire sur 16 bits du nombre qui s'écrit 127 
en base 10$\:$?

\tiret Quelle est l'écriture décimale du nombre entier positif re\-pré\-sen\-té
en binaire sur 16 bits par  0011010101010101$\:$?

\tiret Quel est le texte représenté en code {\sc ascii} par les octets suivants~:
01000100 01100101 01110011 01100011 01100001
01110010 0111\-0100 0110\-0101 0111\-0011$\:$?

\noindent $\Box$}\end{exercices}
\begin{figure}[htbp]
 \includevspacePostScript[9.6cm]{TelegrapheChappe}{0.95}{0.95}
 \caption{Alphabet du télégraphe optique de Chappe}\label{fig:chappe}
\end{figure}
\begin{table}[htbp]
\begin{center}
\begin{tabular}{@{}|l|l||l|l||l|l||l|l||l|l|@{}} \hline
A & Alpha   & B  & Bravo & C & Charlie & D & Delta   & E  & Echo \\\hline
F & Foxtrot  & G & Golf   & H  & Hotel & I & India & J & Juliet \\\hline
K  & Kilo & L & Lima & M & Mike   & N  & November & 0 & Oscar \\\hline
P & Papa   & Q  & Quebec & R & Romeo  & S & Sierra   & T  & Tango \\\hline
U & Uniform & V & Victor   & W  & Whiskey & X & Xray & Y & Yankee   \\\hline
Z  & Zulu &   \multicolumn{8}{c}{}   \\\cline{1-2}
\end{tabular}
\end{center}
\caption{Alphabet phonétique international}\label{tab:alphone}
\end{table}
\begin{exercice}{Les chansons sont des sources inépuisables d'exercices
 \cite{boutet}, \cite{sabatier1}, \cite{sabatier2}, \cite{sabatier3}. 
 Écrire par exemple des programmes 
\pascal\  pour jouer {\em Green\-sleeves\/} (figure \ref{fig:Greensleeves}, page \pageref{fig:Greensleeves})
et {\em O Tannenbaum\/} (figure \ref{fig:Tannenbaum}, page \pageref{fig:Tannenbaum}).
On pourra également dérouler la partition musicale ligne par ligne sur l'écran 
en même temps qu'elle est jouée par l'ordinateur en re\-pré\-sen\-tant la portée
et les notes comme  à la figure \ref{fig:RepNotes} (page \pageref{fig:RepNotes}).
\noindent $\Box$}\end{exercice}
\begin{figure}[htbp]
 %\includevspacePostScript[12.4cm]{Greensleeves}{0.95}{0.95}
 %%&frLaTeX
%\documentclass{article}
%\usepackage{musixtex}
%\begin{document}
\begin{flushleft}
\parindent0pt
\begin{music}
\generalmeter{\meterfrac34}
\nostartrule
\parindent0pt%
\nobarnumbers
\startpiece
%
\NOtes
\Uptext{\metron{\qu}{144} Allegro}%
\qp\qp\qu{h}\enotes\doublebar
\NOtes\hl{j}\ql{k}\enotes\bar
\NOtes\qlp{l}\cl{^m}\ql{l}\enotes\bar
\NOtes\hl{k}\ql{i}\enotes\bar
\NOtes\qup{g}\cu{h}\ql{i}\enotes\bar
%
\NOtes\hl{j}\qu{h}\enotes\bar
\NOtes\qup{h}\cu{^g}\qu{h}\enotes\bar
\NOtes\hl{i}\qu{^g}\enotes\bar
\NOtes\hu{e}\qu{h}\enotes\bar
\NOtes\hl{j}\ql{k}\enotes\bar
\NOtes\qlp{l}\cl{^m}\ql{l}\enotes\bar
%
\NOtes\hl{k}\ql{i}\enotes\bar
\NOtes\qup{g}\cu{h}\qu{i}\enotes\bar
\NOtes\qlp{j}\cl{i}\qu{h}\enotes\bar
\NOtes\qup{^g}\cu{^f}\qu{g}\enotes\bar
\NOtes\hup{h}\enotes\bar
%
\NOtes\qu{h}\ibl013\qb0{h}\qb0{i}\qb0{j}\tbl0\qb0{l}\enotes\leftrepeat
\NOtes\hlp{n}\enotes\bar
\NOtes\qlp{n}\cl{^m}\ql{l}\enotes\bar
\NOtes\hl{k}\ql{i}\enotes\bar
\NOtes\qup{g}\cu{h}\ql{i}\enotes\bar
\NOtes\hl{j}\qu{h}\enotes\bar
%
\NOtes\qup{h}\cu{^g}\qu{h}\enotes\bar
\NOtes\hl{i}\qu{^g}\enotes\bar
\NOtes\hup{e}\enotes\bar
\NOtes\hup{n}\enotes\bar
\NOtes\qlp{n}\cl{^m}\ql{l}\enotes\bar
\NOtes\hl{k}\ql{i}\enotes\bar
%
\NOtes\qup{g}\cu{h}\ql{i}\enotes\bar
\NOtes\qlp{j}\cl{i}\qu{h}\enotes\bar
\NOtes\qup{^g}\cu{^f}\qu{g}\enotes\bar
\NOtes\hup{h}\enotes\Setvolta1\bar
\NOtes\qu{h}\ibl013\qb0{h}\qb0{i}\qb0{j}\tbl0\qb0{l}\enotes\setendvoltabox\Setvolta2\rightrepeat%
\NOtes\hu{h}\qp\enotes\setendvoltabox
%
\Endpiece
\end{music}
\end{flushleft}%
%\end{document}%
\endinput
 \caption{Greensleeves (chanson irlandaise)}\label{fig:Greensleeves}
\end{figure}
\begin{figure}[htbp]
 %\includevspacePostScript[2.8cm]{OTannenbaum}{0.95}{0.95}
 %%&frLaTeX
%\documentclass{article}
%\usepackage{musixtex}
%\begin{document}
\begin{flushleft}
\parindent0pt
\begin{music}
\generalmeter{\meterfrac34}
\nostartrule
\parindent0pt%
\nobarnumbers
\startpiece
%
\notes\qp\qp\qu{c}\enotes\bar
\notes\cup{f}\ccu{f}\qu{f}\qu{g}\enotes\bar
\notes\cup{h}\ccu{h}\qu{h}\qu{f}\enotes\bar
\notes\cu{g}\cu{h}\ql{i}\qu{e}\enotes\bar
\notes\qu{g}\qu{f}\enotes\setrightrepeat\alaligne
%
\notes\qp\qp\ds\cl{j}\enotes\bar
\notes\cl{j}\cu{h}\qlp{k}\cl{j}\enotes\bar
\notes\cl{j}\cl{i}\qlp{i}\cl{i}\enotes\bar
\notes\cl{i}\cl{g}\qlp{j}\cl{i}\enotes\bar
\notes\cl{i}\cu{h}\qu{h}\enotes\setrightrepeat
%
\endpiece
\end{music}
\end{flushleft}%
%\end{document}%
\endinput
 \caption{O Tannenbaum  (chanson allemande)}\label{fig:Tannenbaum}
\end{figure}
\begin{figure}[htbp]
 \includevspacePostScript[2.3cm]{EcrireEtJouerMus}{0.95}{0.95}
 \caption{Portée et notes représentées sur l'écran de l'ordinateur}\label{fig:RepNotes}
\end{figure}
\begin{corrige}{
Nous donnons un programme, les autres se trouvant sur disquette~:
\input CodeMorse1.tex

On trouvera sur la disquette diverses améliorations possibles de ce programme, notamment
un retour automatique à la ligne.
\noindent $\Box$}\end{corrige}
\begin{corriges}{\ \\
\vspace{-\baselineskip}

\tiret La représentation binaire sur 16 bits de 127 est 
00\-00\-00\-00\-01\-11\-11\-11.

\tiret Le nombre représenté par le demi-mot 0011010101010101 s'écrit
13653 en base 10.

\tiret La suite d'octets 01000100 01100101 01110011 01\-10\-00\-11 01\-10\-00\-01
01\-11\-00\-10 01110100 01100101 01110011
soit, en notation décimale 68 101 115 99 97 114 116 101 115,
représente la chaîne de caractères
{\cha $\scriptstyle '$Descartes\idt $\scriptstyle '$}.

\noindent $\Box$}\end{corriges}
\begin{corrige}{
Nous donnons un programme, de nombreuses autres chansons se trouvant programmées sur disquette~:
\input OTannenbaum.tex
\noindent $\Box$}\end{corrige}

%%%%%%%%%%%%%%%%%%%%%%%%%%%%%%%%%%%%%%%%%%%%%%%%%%%%%%%%%%%%%%%%%%%%%%%%%%%
\noindent \framebox[\textwidth]{
\begin{minipage}{\myparboxwidth}
\setlength{\parindent}{2em}
\small \sf
\vspace{\abovedisplayskip}
\section*{Addenda}

En complément nous expliquons deux représentations des entiers relatifs, la se\-con\-de
étant la plus employée. 

\paragraph*{Représentation signée des entiers relatifs~:\ }
\index{Numération binaire}
Pour coder un entier relatif (positif ou négatif) sur $n$ bits, on peut réserver un bit 
pour le signe et coder la valeur absolue de l'entier sur les $n - 1$ bits restants. 
Prenons la convention
de placer le signe en tête en choisissant 0 pour le signe $+$ et 1 pour le signe $-$. 
Le code de
$-10$ sur un octet sera 10001010 soit 1 pour le signe suivi de 0001010 pour la valeur absolue $10$ 
de $-10$. Avec
cette convention on peut coder sur un octet les entiers compris entre $-127$ (11111111) et $127$
(01111111). L'entier nul est représenté par 00000000 ($+0$) ou bien par 10000000 ($-0$).

\vspace{\abovedisplayskip}
\end{minipage}
}
%%%%%%%%%%%%%%%%%%%%%%%%%%%%%%%%%%%%%%%%%%%%%%%%%%%%%%%%%%%%%%%%%%%%%%%%%%%% 

%%%%%%%%%%%%%%%%%%%%%%%%%%%%%%%%%%%%%%%%%%%%%%%%%%%%%%%%%%%%%%%%%%%%%%%%%%%
\noindent \framebox[\textwidth]{
\begin{minipage}{\myparboxwidth}
\setlength{\parindent}{2em}
\small \sf
\vspace{\abovedisplayskip}
\paragraph*{Représentation en complément à deux des entiers relatifs~:\ }
Dans la représentation en com\-plé\-ment à deux sur $n$ bits, les nombres positifs 
strictement inférieurs à  $2^{n-1}$ sont
re\-pré\-sen\-tés comme précédemment en base 2 sur $n-1$ bits, le premier bit de signe étant 
nul. Un nombre négatif $-x$
tel que $x \leq 2^{n-1}$ est re\-pré\-sen\-té par un premier bit de signe égal à 1 suivi de la
représentation du nombre positif $p < 2^{n-1}$ tel que $-x = -2^{n-1} + p$
codé
sur  $n - 1$ bits. Par exemple pour coder - 1 sur
un octet, on écrit $-1 = -128 + 127$ et l'on code 127 sur 7 bits ce qui donne 1111111 que l'on fait
précéder de 1 pour trouver 11111111. Avec cette convention on peut également coder
$-128 = -128 + 0$ ce qui donne le code 10000000. Cette méthode permet de 
coder sur un octet tous les entiers relatifs compris entre -128 et +127. \index{Entier}Sur 16 bits, on pourra
coder les entiers compris entre -32768 et 32767 ce qui correspond le plus souvent aux variables 
entières de type {\idt integer} en \pascal. Sur 32 bits on pourra coder les entiers compris
entre $-2\,147\,483\,648$ et $2\,147\,483\,647$ ce qui correspond le plus souvent aux va\-ria\-bles 
entières de type {\idt longint} en \pascal. Cette représentation conduit à des algorithmes
d'addition et soustraction très simples. De plus l'algorithme pour le codage d'un nombre négatif
$-x$ est également simple. On code le nombre positif $x$ sur $n$ bits puis on inverse les bits 
0 en 1 et les bits 1 en 0 puis on ajoute 1. Par exemple pour coder -127, on code 127 en 01111111 puis on
inverse tous les bits ce qui donne 10000000 et on ajoute 1 pour trouver 10000001.
\vspace{\abovedisplayskip}
\end{minipage}
}
%%%%%%%%%%%%%%%%%%%%%%%%%%%%%%%%%%%%%%%%%%%%%%%%%%%%%%%%%%%%%%%%%%%%%%%%%%%% 
\newpage                 % Si on termine le chapitre par une page blanche
\thispagestyle{empty}

